\subsubsection{把一组原子移到目标中心}

居中时，同一组中的所有原子使用同一个位移。设选择中有 $N$ 个原子，分子用笛卡尔位置 $\mathbf r_i$，晶体用分数坐标 $\mathbf f_i$，两种情况分别为
\begin{equation}
 \begin{aligned}
 \Delta\mathbf r&=\mathbf r_{\mathrm{target}}-
        \frac1N\sum_{i=1}^{N}\mathbf r_i,
 &\mathbf r_i'&=\mathbf r_i+\Delta\mathbf r,\\
 \Delta\mathbf f&=\mathbf f_{\mathrm{target}}-
        \frac1N\sum_{i=1}^{N}\mathbf f_i,
 &\mathbf f_i'&=\mathbf f_i+\Delta\mathbf f.
 \end{aligned}
 \label{eq:center-selected-atoms}
\end{equation}
第一行用于分子，位置与位移的单位都是\AA；第二行用于晶体，各量均为无量纲分数坐标。分子常用目标 $(0,0,0)$，晶体常用 $(0.5,0.5,0.5)$。每个原子在平均值中占相同权重，因此居中不使用原子质量，也不改变所选原子之间的距离。

例如，分子中两个原子的 $x$ 坐标为1和3\,\AA，要把它们的平均位置移到原点，共同位移就是 $-2\,\text{\AA}$；处理后为 $-1$ 和1\,\AA，间距仍为2\,\AA。晶体中若两个位点的第一分数分量为0.1和0.3，目标为0.5，则都加0.3，得到0.4和0.6。这里只取一个分量说明计算，其余两个分量同样处理。

周期边界附近要先看清使用的是哪一个映像。0.99与0.01的普通平均为0.5，但两个位点可能实际聚集在边界两侧；若要以这段连续结构的中心为准，应先整理为相邻的连续映像再居中。应用后坐标可能被归回晶胞，因此还应从周期连接看整体，而不只比较某个坐标的正负。

\subsubsection{绕指定锚点旋转}

旋转先从每个位置中减去锚点，再转动这个相对向量，最后加回锚点。这样，锚点本身不动，离锚点的距离也不变。沿用行向量表示，有
\begin{equation}
 \mathbf r'=\mathbf p+(\mathbf r-\mathbf p)
              R(\hat{\mathbf n},\theta)^{\mathsf T}.
 \label{eq:rotate-row-position}
\end{equation}
$\mathbf p$ 是以\AA 为单位的锚点，$\hat{\mathbf n}$ 是旋转轴的单位方向，$\theta$ 是转角；上标 $\mathsf T$ 表示转置。对话框中填写度数，正角遵循右手定则：右手拇指沿输入轴的正向，四指弯曲的方向为正转向。轴向量只决定方向，例如\literalinput{0 0 1}和\literalinput{0 0 2}表示同一根正向转轴。

绕全局 $z$ 轴旋转时，常用的矩阵为
\[
 R_z(\theta)=\begin{bmatrix}
 \cos\theta&-\sin\theta&0\\
 \sin\theta&\cos\theta&0\\
 0&0&1
 \end{bmatrix}.
\]
取锚点 $\mathbf p=(1,1,0)\,\text{\AA}$，位置 $\mathbf r=(2,1,0)\,\text{\AA}$，绕 $z$ 轴正转 $90^\circ$，得到 $\mathbf r'=(1,2,0)\,\text{\AA}$。这个小例子可用来辨认旋转方向；旋转后原子到锚点仍相距1\,\AA。晶体中原子的等价位置还会归回晶胞，预览时应连同边界映像一起观察。

\subsubsection{关于一个平面作镜像}

镜像需要一个平面上的点 $\mathbf p$ 和非零笛卡尔法向 $\mathbf n$。先把法向归一化，再减去两倍的有向投影：
\begin{equation}
 \begin{aligned}
 \hat{\mathbf n}&=\frac{\mathbf n}{\|\mathbf n\|},\\
 \mathbf r'&=\mathbf r-
 2\bigl[(\mathbf r-\mathbf p)\cdot\hat{\mathbf n}\bigr]
       \hat{\mathbf n}.
 \end{aligned}
 \label{eq:mirror-cartesian-plane}
\end{equation}
点积 $(\mathbf r-\mathbf p)\cdot\hat{\mathbf n}$ 是原子到平面的有向距离，单位为\AA。镜像后它的符号反转，大小相同；位于平面上的原子保持不动。法向填\literalinput{1 0 0}，平面点取 $(1,0,0)\,\text{\AA}$，就是 $x=1\,\text{\AA}$ 的平面。位置 $(1.5,2,0)\,\text{\AA}$ 镜像后变为 $(0.5,2,0)\,\text{\AA}$，两侧到平面均为0.5\,\AA。

\begin{figure}[!htb]
\centering
\begin{tikzpicture}[x=1mm,y=1mm,font=\sffamily\small]
 \node[text=brand] at (24,47) {绕锚点旋转};
 \draw[teal,dashed] (10,15)--(33,15);
 \draw[teal,dashed] (10,15)--(10,38);
 \draw[->,>=stealth,brand,line width=.85pt] (33,15) arc[start angle=0,end angle=90,radius=23];
 \fill[ink] (10,15) circle (1.1);
 \fill[teal] (33,15) circle (1.4);
 \fill[brand] (10,38) circle (1.4);
 \node[anchor=north] at (10,13) {$\mathbf p$};
 \node[anchor=west,text=teal] at (35,15) {$\mathbf r$};
 \node[anchor=east,text=brand] at (8,38) {$\mathbf r'$};
 \node[text=brand] at (30,35) {$+90^\circ$};
 \node[text=muted,font=\sffamily\footnotesize] at (24,3) {到锚点的距离保持不变};
 \node[text=brand] at (86,47) {关于平面镜像};
 \fill[paleblue] (83,8) rectangle (87,40);
 \draw[teal,line width=.8pt] (85,8)--(85,40);
 \draw[dashed,muted] (68,29)--(102,29);
 \draw[->,>=stealth,brand,line width=.7pt] (85,12)--(105,12) node[above] {$\hat{\mathbf n}$};
 \fill[ink] (85,18) circle (1.1);
 \fill[teal] (102,29) circle (1.4);
 \fill[brand] (68,29) circle (1.4);
 \node[anchor=south,text=teal] at (102,31) {$\mathbf r$};
 \node[anchor=south,text=brand] at (68,31) {$\mathbf r'$};
 \node[anchor=east] at (83,18) {$\mathbf p$};
 \node[text=muted,font=\sffamily\footnotesize] at (86,3) {到镜面的距离两侧相等};
\end{tikzpicture}
\caption{锚点旋转与平面镜像（截面示意）。左图中的点绕 $\mathbf p$ 旋转；右图中的点跨过镜面，沿法向到镜面的距离大小不变。}
\label{fig:round10-anchor-mirror}
\end{figure}


上述三个例子只说明参数的几何关系，不代表一次能量优化。实际处理分子时，先确认所选原子构成需要整体移动的部分，再按下面的步骤选择一种操作；不要把相机转动当作结构已经旋转。
